\documentclass[]{beamer}

\usepackage[utf8]{inputenc}
\usepackage[portuges]{babel}
\usepackage[T1]{fontenc}
\usepackage{hyperref}
\usepackage[]{csquotes}

\setbeamertemplate{footline}[page number]
\setbeamercovered{transparent}
\beamertemplatenavigationsymbolsempty

\def\BLACKBG{true}
\ifdefined\BLACKBG
  \setbeamercolor{block body}{bg=normal text.bg!90!black}
  \setbeamercolor{block body example}{bg=normal text.bg!90!black}
  \setbeamercolor{block title alerted}{use={normal text,alerted text},fg=alerted text.fg!75!normal text.fg,bg=normal text.bg!75!white}
  \setbeamercolor{block title}{bg=black}
  \setbeamercolor{block title example}{use={normal text,example text},fg=example text.fg!75!normal text.fg,bg=normal text.bg!75!white}
  \setbeamercolor{fine separation line}{}
  \setbeamercolor{frametitle}{fg=white}
  \setbeamercolor{item projected}{fg=white}
  \setbeamercolor{normal text}{bg=black,fg=white}
  \setbeamercolor{separation line}{}
  \setbeamercolor{structure}{bg=black,fg=white}
  \setbeamercolor{title}{fg=white}
  \setbeamercolor{titlelike}{fg=white}
\else
  \setbeamercolor{block body}{bg=normal text.bg!90!white}
  \setbeamercolor{block body example}{bg=normal text.bg!90!white}
  \setbeamercolor{block title alerted}{use={normal text,alerted text},fg=alerted text.fg!75!normal text.fg,bg=normal text.bg!75!black}
  \setbeamercolor{block title}{bg=white}
  \setbeamercolor{block title example}{use={normal text,example text},fg=example text.fg!75!normal text.fg,bg=normal text.bg!75!black}
  \setbeamercolor{fine separation line}{}
  \setbeamercolor{frametitle}{fg=black}
  \setbeamercolor{item projected}{fg=black}
  \setbeamercolor{normal text}{bg=white,fg=black}
  \setbeamercolor{separation line}{}
  \setbeamercolor{structure}{bg=white,fg=black}
  \setbeamercolor{title}{fg=black}
  \setbeamercolor{titlelike}{fg=black}
\fi

\newcommand{\invertbgcolor}{
  \ifdefined\BLACKBG
    %\setbeamercolor{normal text}{bg=white,fg=black}
  \fi
}

\newcommand{\backbgcolor}{
  \ifdefined\BLACKBG
    %\setbeamercolor{normal text}{bg=white,fg=black}
  \fi
}


\author{Dmitry "Nix" Rocha}
\title{Aprendendo Python com testes}
\institute{ /(guru|pug-)pi/i }

\begin{document}
  \begin{frame}{}
    \titlepage
  \end{frame}

  \begin{frame}{su - mario}
    \tableofcontents
  \end{frame}

  \section{O que é teste automatizado de software?}

  \begin{frame}
    \begin{center}
      \Huge Definições
    \end{center}
  \end{frame}

  \subsection*{Definição - 1}
  \begin{frame}
    \begin{block}{Teste}
      É algo que testo =/
    \end{block}
    \pause

    \begin{block}{Automatizado}
      É algo que automatizo :P
    \end{block}
    \pause

    \begin{block}{Software}
      Conjunto de instruções.
    \end{block}
    \pause
  \end{frame}

  \invertbgcolor
  \begin{frame}
    \begin{center}
      \includegraphics[width=5cm]{images/pokerface}
    \end{center}
  \end{frame}

  \subsection*{Definição - 2}
  \backbgcolor
  \begin{frame}
    \blockquote[{\cite[Wikipedia]{_test_2014}}]{
      In software testing, test automation is the use of special software
      (separate from the software being tested) to control the execution of
      tests and the comparison of actual outcomes with predicted outcomes. Test
      automation can automate some repetitive but necessary tasks in a
      formalized testing process already in place, or add additional testing
      that would be difficult to perform manually.
    }
  \end{frame}

  \invertbgcolor
  \begin{frame}{}
    \begin{center}
      \includegraphics[width=5cm]{images/likeasir}
    \end{center}
  \end{frame}

  \backbgcolor
  \section{Exemplos e casos de uso com Python}

  \begin{frame}{fibonacci}
    Sequência de Fibonacci :)
  \end{frame}

  \bibliographystyle{plain}
  \bibliography{Items}
\end{document}
